\include{../header}

\usepackage{parskip,scalerel}

\newcommand{\diff}[1]{\frac{d}{d#1}}

\begin{document}
\section{PS6: Identifying structure in derivative problems}

So exactly the hard part about learning derivative shortcut rules is figuring out which ones to use when. The direct answer to this question is, you use the shortcut rules that match the \textbf{algebraic structure} of the function whose derivative you are finding.

Identifying \textbf{algebraic structure} takes some practice! Let's get into it with this problem set.

You will need six
different color pens, pencils, highlighters, or whatever, and I recommend printing this sheet out and writing on it directly.

\hrulefill

First, here is a reference list of all the derivative shortcut rules we know so far. (We'll learn more shortcut rules after spring break!)
{ \everymath{\displaystyle}
    \begin{itemize}[label=\(\square\)]
        \item sum rule: 
        \[\diff{x}[f(x)+g(x)] = \diff{x}[f(x)]+\diff{x}[g(x)] \]
        \item constant multiple rule: if $c$ is a constant, then
        \[ \diff{x}[c\cdot f(x)] = c\cdot\diff{x}[f(x)] \]
        \item constant function rule: if $c$ is a constant, then
        \[\diff{x}[c] = 0\]
        \item power rule: 
        \[\diff{x}[x^n] = nx^{n-1}\]
        \item exponential function rule: 
        \[\diff{x}[a^x] = a^x\cdot\ln(a)\]
        \item[] trigonometric function rules: 
        \begin{itemize}[label=\(\bigcirc\)]
            \item $\diff{x}[\sin(x)] = \cos(x)$
            \item $\diff{x}[\cos(x)] = -\sin(x)$
        \end{itemize}
        \item product rule (aka rule for functions that are multiplied): 
        \[\diff{x}[f(x) \cdot g(x)] = \diff{x}[f(x)] \cdot g(x) + \diff{x}[g(x)] \cdot f(x)\]
        \item quotient rule (aka rule for functions that are divided): 
        \[\diff{x}\left[\frac{f(x)}{g(x)}\right] 
        = \frac{g(x) \cdot \diff{x}[f(x)] - f(x) \cdot \diff{x}[g(x)]}{[g(x)]^2}\]
    \end{itemize}
}

\begin{enumerate}[leftmargin=15pt]
    \item Consider the function 
    \[\scaleto{f(t) = \frac{7}{t} - 7^t + 3t^7 - 2\, \sin(t)+ \sqrt{t}}{40 pt} \]
    In different colors for each:
    \begin{enumerate}[(a), leftmargin=15pt]
        \item circle anywhere that functions are being \textit{added or subtracted} to make a bigger function,
        \item circle everything that's a \textit{constant multiple} of some function,
        \item circle everything that's a \textit{power function} (hint: there are three!),
        \item circle everything that's an \textit{exponential function},
        \item and lastly circle everything that's a \textit{trigonometric function}.
    \end{enumerate}
    Finally, now that you have done all this circling, find $\displaystyle\diff{t} [f(t)]$ using the appropriate shortcut rules. The more work you can show, the more confidence you will gain in the process!
    \pagebreak

    \renewcommand*{\thefootnote}{\fnsymbol{footnote}}

    \item Consider the function\footnote{I pulled this function from \href{https://library.tbil.org/2025/calculus/exercises/\#/bank/DF4/1/}{CheckIt practice problems DF4}!} 
    \[\scaleto{g(x) = (2x^2 - 5x - 3) \cos(x)}{30 pt} \]
    In different colors for each:
    \begin{enumerate}[(a), leftmargin=15pt]
        \item circle anywhere that functions are being \textit{multiplied} to make a bigger function,
        \item circle anywhere that functions are being \textit{added or subtracted} to make a bigger function,
        \item circle everything that's a \textit{constant multiple} of some function,
        \item circle everything that's a \textit{power function} (hint: there are two!),
        \item circle everything that's a \textit{constant function},
        \item and lastly circle everything that's a \textit{trigonometric function}.
    \end{enumerate}
    Finally, now that you have done all this circling, find $\displaystyle\diff{x} [g(x)]$ using the appropriate shortcut rules. The more work you can show, the more confidence you will gain in the process!
    \vfill

    \pagebreak

    \item Consider the function\footnote{This function is also from \href{https://library.tbil.org/2025/calculus/exercises/\#/bank/DF4/1/}{CheckIt practice problems DF4}!} 
    \[\scaleto{h(y) = \dfrac{e^y}{3y^2+3y-5}}{50 pt} \]
    In different colors for each:
    \begin{enumerate}[(a), leftmargin=15pt]
        \item circle anywhere that functions are being \textit{divided} to make a bigger function,
        \item circle anywhere that functions are being \textit{added or subtracted} to make a bigger function,
        \item circle everything that's a \textit{constant multiple} of some function,
        \item circle everything that's a \textit{power function} (hint: there are two!),
        \item circle everything that's a \textit{constant function},
        \item and lastly circle everything that's an \textit{exponential function}.
    \end{enumerate}
    Finally, now that you have done all this circling, find $\displaystyle\diff{y} [h(y)]$ using the appropriate shortcut rules. The more work you can show, the more confidence you will gain in the process!
    \vfill

\end{enumerate}


\end{document}